
\usepackage[left=1cm,right=1cm,top=1.5cm,bottom=2cm]{geometry}
\usepackage{xcolor}
\usepackage{array}
\newcolumntype{R}[1]{>{\raggedleft\let\newline\\\arraybackslash\hspace{0pt}}p{#1}}
  % \usepackage[finnish]{babel}
  % % Modifying the title
  % 
  \makeatletter
  \renewcommand{\maketitle}{\bgroup\setlength{\parindent}{0pt}
    \begin{flushleft}
    % \LARGE \@title
    
    \vspace{-45mm}
    \Large \@author
    
    \vspace{-3mm}
    \normalsize \@date
    
    \tableofcontents
    
    \end{flushleft}\egroup
  }
  
  
  % ------------------------------------------------ %
    %                                                  %
    %              Fonts and typography                %
    %                                                  %
    % ------------------------------------------------ %
    
    %% Typography
  % \usepackage[no-math]{fontspec}
  \usepackage{color}
  
  %% Fonts
  % \setmainfont{PT Sans}
  % \setsansfont{PT Sans}
  % \setmonofont{PT Sans}
  % 
  % \newfontfamily\PT SansBlack{PT Sans Black}
  
  \usepackage{lmodern}
  % \usepackage{helvet} % or 
  % \usepackage{tgheros}
% 
%%% we want that the default font family is sans serif
\renewcommand{\rmdefault}{\sfdefault}
  
\usepackage{etoolbox}
\AtBeginEnvironment{tabular}{\sffamily}
  
  %% Sans serif as main font
  %\renewcommand{\familydefault}{\sfdefault}
  
  %% Set Sans font in headings
  %\usepackage{sectsty}
  %\allsectionsfont{\sffamily}
  
  %% Set polyglossia language
  % \usepackage{polyglossia}
  % \setdefaultlanguage{finnish}
  % disable hyphenation
  %\hyphenpenalty = 10000
  
  %% Emphasizing inline code output
  \usepackage{alltt}
  
  
  % ------------------------------------------------ %
    %                                                  %
    %              Settings for knitr                  %
    %                                                  %
    % ------------------------------------------------ %
    
    %% for Shaded environment in knitr
  \usepackage{framed}
  \usepackage{fancyhdr}
  
  %% landscape mode
  \usepackage{lscape}
  
  %% Encoding
  %\usepackage[UTF-8]{inputenc}
  
  % rowspacing
  % \usepackage{setspace}
  
  % ------------------------------------------------ %
    %                                                  %
    %              Page settings                       %
    %                                                  %
    % ------------------------------------------------ %
    
    % \input{preamb_a4.tex}
  
  
  %% Set line spacing
  \usepackage{setspace}
  \setstretch{1.0}
  
  %% Disable paragraph indentation
  \usepackage{parskip}
  
  %% Custom headers
  % \usepackage{fancyhdr}
  
  %% Multicolumns
  \usepackage{multicol}
  
  % ------------------------------------------------ %
    %                                                  %
    %                  Headings                        %
    %                                                  %
    % ------------------------------------------------ %
    
    
    \usepackage{enumerate}
  
  
  % ------------------------------------------------ %
    %                                                  %
    %               Table of content                   %
    %                                                  %
    % ------------------------------------------------ %
    
    % dotted lines for table of content
  
  \makeatletter
  \renewcommand*\l@section{\@dottedtocline{1}{1.5em}{2.3em}}
  \makeatother
  \setcounter{tocdepth}{1}
  
  % ------------------------------------------------ %
    %                                                  %
    %               % Colors                           %
  %                                                  %
    % ------------------------------------------------ %
    
    
    %% Tango color scheme
  \definecolor{SkyBlue}{HTML}{3465A4}
  \definecolor{DarkSkyBlue}{HTML}{204A87}
  \definecolor{Plum}{HTML}{75507B}
  \definecolor{ScarletRed}{HTML}{CC0000}
  \definecolor{Aluminium1}{HTML}{EEEEEC}
  \definecolor{Aluminium6}{HTML}{2e3436}
  \definecolor{Black}{HTML}{000000}
  
  % ------------------------------------------------ %
    %                                                  %
    %               Links                              %
    %                                                  %
    % ------------------------------------------------ %
    
    %% Hyperref
  \usepackage{hyperref}
  
  \hypersetup {
    linkcolor = Aluminium6,
    citecolor = DarkSkyBlue,
    filecolor = DarkSkyBlue,
    urlcolor = DarkSkyBlue
  }
  
  %% Don’t use Mono font for URLs
  \urlstyle{same}
  
  
  % ------------------------------------------------ %
    %                                                  %
    %               Images & tables                    %
    %                                                  %
    % ------------------------------------------------ %
    
    % Images
  \usepackage{graphicx}
  % tables
  \usepackage{longtable}
  \usepackage{booktabs}
  \usepackage{float}
  \restylefloat{table}
  \usepackage{wrapfig}
  
  \renewcommand\arraystretch{1.1} % table rowheight
  % see: http://everythingyouforgetaboutlatex.blogspot.fi/2008/10/formatting-tables.html
  \setlength{\intextsep}{-0.1cm} % spacing between floats
  
  
  % SHADOWS FOR FLAGS
  
  \usepackage{tikz}
  \usetikzlibrary{shadows,calc}
  
  % code adapted from https://tex.stackexchange.com/a/11483/3954
  
  % some parameters for customization
  \def\shadowshift{3pt,-3pt}
  \def\shadowradius{6pt}
  
  \colorlet{innercolor}{black!60}
  \colorlet{outercolor}{gray!05}
  
  % this draws a shadow under a rectangle node
  \newcommand\drawshadow[1]{
    \begin{pgfonlayer}{shadow}
    \shade[outercolor,inner color=innercolor,outer color=outercolor] ($(#1.south west)+(\shadowshift)+(\shadowradius/2,\shadowradius/2)$) circle (\shadowradius);
      \shade[outercolor,inner color=innercolor,outer color=outercolor] ($(#1.north west)+(\shadowshift)+(\shadowradius/2,-\shadowradius/2)$) circle (\shadowradius);
        \shade[outercolor,inner color=innercolor,outer color=outercolor] ($(#1.south east)+(\shadowshift)+(-\shadowradius/2,\shadowradius/2)$) circle (\shadowradius);
          \shade[outercolor,inner color=innercolor,outer color=outercolor] ($(#1.north east)+(\shadowshift)+(-\shadowradius/2,-\shadowradius/2)$) circle (\shadowradius);
            \shade[top color=innercolor,bottom color=outercolor] ($(#1.south west)+(\shadowshift)+(\shadowradius/2,-\shadowradius/2)$) rectangle ($(#1.south east)+(\shadowshift)+(-\shadowradius/2,\shadowradius/2)$);
              \shade[left color=innercolor,right color=outercolor] ($(#1.south east)+(\shadowshift)+(-\shadowradius/2,\shadowradius/2)$) rectangle ($(#1.north east)+(\shadowshift)+(\shadowradius/2,-\shadowradius/2)$);
                \shade[bottom color=innercolor,top color=outercolor] ($(#1.north west)+(\shadowshift)+(\shadowradius/2,-\shadowradius/2)$) rectangle ($(#1.north east)+(\shadowshift)+(-\shadowradius/2,\shadowradius/2)$);
                  \shade[outercolor,right color=innercolor,left color=outercolor] ($(#1.south west)+(\shadowshift)+(-\shadowradius/2,\shadowradius/2)$) rectangle ($(#1.north west)+(\shadowshift)+(\shadowradius/2,-\shadowradius/2)$);
                    \filldraw ($(#1.south west)+(\shadowshift)+(\shadowradius/2,\shadowradius/2)$) rectangle ($(#1.north east)+(\shadowshift)-(\shadowradius/2,\shadowradius/2)$);
                      \end{pgfonlayer}
  }
  
  % create a shadow layer, so that we don't need to worry about overdrawing other things
% \pgfdeclarelayer{shadow} 
% \pgfsetlayers{shadow,main}
% 
% \newsavebox\mybox
% \newlength\mylen
% 
% \newcommand\shadowimage[2][]{%
% \setbox0=\hbox{\includegraphics[#1]{#2}}
% \setlength\mylen{\wd0}
% \ifnum\mylen<\ht0
% \setlength\mylen{\ht0}
% \fi
% \divide \mylen by 120
% \def\shadowshift{\mylen,-\mylen}
% \def\shadowradius{\the\dimexpr\mylen+\mylen+\mylen\relax}
% \begin{tikzpicture}
% \node[anchor=south west,inner sep=0] (image) at (0,0) {\includegraphics[#1]{#2}};
% \drawshadow{image}
% \end{tikzpicture}}




% ------------------------------------------------ %
%                                                  %
%               First page graphics                %
%                                                  %
% ------------------------------------------------ %

%\usepackage[showframe,textwidth=3.0in]{geometry}
\usepackage{xstring}% for string comparrison
\usepackage{calc}%    for \widthof
\usepackage{pgf}%     for math calclations


\usepackage{caption}
\captionsetup{labelformat=empty, justification=justified, font={large,it}}
\usepackage{float}
\floatstyle{plaintop}
\restylefloat{table}

% https://texfaq.org/FAQ-vertposfp
\makeatletter
\setlength{\@fptop}{0pt}
\makeatother


% <!--- Define Headers and Footers --->
% \fancypagestyle{plain}{%
%   \renewcommand{\headrulewidth}{1pt}%
%   \fancyhf{}%
%   \setlength\footskip{0pt}
% }
% \pagestyle{plain}
% %
% \setlength\headheight{35pt}
% \lhead{Huono-osaisuus Suomessa - Alueprofiilit}
% \rhead{\includegraphics[width=2.0cm]{diak-logo-1.pdf}}
% \fancyfoot[C]{\thepage}

% \pagestyle{fancy}
% \setlength\headheight{35pt}
% \fancyhf{}
% \lhead{\Large{Huono-osaisuus Suomessa - Alueprofiilit}}
% \rhead{\includegraphics[width=1.0cm]{diak-logo-1.pdf}}
% \fancyfoot[C]{\thepage}

% \pagestyle{fancy}
% \head{}
% \cfoot{}

\thispagestyle{empty}

